% Pursuing Stacks (cotes 134-2 to 134-8): archive cote 134-2, site page 18 (PDF page 19).
% letter to Quillen. Draft compared: ps1.tex:216-293.
% AI-transcribed by Claude Opus 5.5 (claude-opus-5-5, adaptive thinking, effort high), Message Batches API, batch ps1 (msgbatch_01Rv5hCaNTtUt4JHx1yqoHPL), 2026-11-12; not yet checked by eye.
\scanpage{134-2/18}
\margin{3)}

\mspage{5}

concerning this secondary structure. The first thing I could think of is
the ``pentagon axiom'' for the associativity data, which occurs when
\margin{$\overset{i}{\underset{\ell}{*}}$} looking at associativities for the compositum (for $\overset{\hand{i}}{\underset{\hand{\ell}}{*}}$ say) of four factors.
Here again\xout{x} the first reflex would be to write down, as usual, an
\uline{equality} for \xout{the} two compositions of associativity isomorphisms, exhibited
in the pentagon diagram. One suspects however that such equality should,
again, be replaced by a ``homotopy''-arrow , which now appears as a kind
of ``ternary'' structure - before even having exhausted the list of coheren%
ce\hand{``}relations\hand{''}one could think of with respect to the secondary structure !
Here one seems caught at first sight in an infinite chain of \deleted{increasingly}\inserted{ever}
\hand{``}higher\hand{''}, and presumably\hand{,} messier structures, where one is going to get hope%
-lessly lost, unless one discovers some simple guiding principle for she\over{c}{d}-%
ding some clarity in\hand{to}\hand{|}the mess.

\margin{\unsure{3.}} \hand{3.} I thought of writing you mainly because I believe that, if anybody,
you should know if the kind of structure I am looking for has been worked
out - maybe even \uline{you} did ? In this respect, I vaguely remember that
you had a description of n-categories in terms of n-semisimplicial sets,
satisfying certain exactness conditions, in much the same way as an ordina-%
ry category can be interpreted, via its ``nerve'', as a particular type of
semisimplicial set. But I have no idea if your definition applied only
for describing n-categories with strict associativities, or not\hand{$^{(6)}$}\hand{.}

Still some comments in the spirit of \deleted{your}\inserted{Quillen's} axiomatics of homotopical
algebra - in order to make the question I am proposing more seducing
maybe to you ! One comment is that presumably, the category of $\infty$-grou-%
poids (which is still to be defined) is a ``model category'' for the usual
homotopy category; this would be at any rate one plausible way to make
explicit\hand{\unsure{|}}the intuition referred to before, that a homotopy type is ``essen%
tially the same'' as an $\infty$-groupoid up to $\infty$-equivalence. The other
comment: the construction of the fundamental $\infty$-groupoid of a space,
disregarding for the time being the question of working out in full the
pertinent structure on this messy object, can be paraphrased in any model
category in your sense, and yields a functor from this category to the
category of $\infty$ - groupoids, and hence (by geometric realization, or
by localization) also to the usual homotopy category\hand{$^{(7)}$}. Was this functor
obvious beforehand ? It is of a non-trivial nature only when the model
category is \uline{not} pointed - as a matter of fact the whole construction
can be carried out canonically, in terms of a ``cylinder object'' I for
the final object e of the model category\hand{,} \hand{playing the role of the unit segment.}

\margin{18}
